\documentclass[../../../include/open-logic-section]{subfiles}

\begin{document}

\olfileid[id]{sth}{ordinals}{vn}
\olsection[Konstruksi von Neumann]{Konstruksi von Neumann atas Ordinal}

\olref[sth][ordinals][iso]{thm:woalwayscomparable} memunculkan suatu gagasan.
Kita dapat memperkenalkan objek-objek tertentu, yang disebut \emph{tipe
urutan}, untuk bertindak sebagai wakil pengurutan baik. Dengan menulis
$\ordtype{A, <}$ untuk tipe urutan dari pengurutan baik $\tuple{A, <}$, kita
berharap dapat memperoleh dua prinsip berikut:
\begin{align*}
	\ordtype{A, <} = \ordtype{B, \lessdot} &
	\text{ jika dan hanya jika } \ordeq{\tuple{A, <}}{\tuple{B, \lessdot}}\\
	\ordtype{A, <} < \ordtype{B, \lessdot}&
	\text{ jika dan hanya jika }\ordeq{\tuple{A, <}}{\tuple{B_b, \lessdot_b}}\text{ untuk suatu }b \in B
\end{align*}
Selain itu, kita mungkin berharap dapat memperkenalkan tipe-tipe urutan
\emph{sebagai himpunan-himpunan tertentu}, sama seperti kita dapat
memperkenalkan bilangan asli sebagai himpunan-himpunan tertentu.

Cara paling umum untuk melakukannya---dan pendekatan yang akan kita
ikuti---adalah mendefinisikan tipe-tipe urutan ini melalui
himpunan-himpunan terurut baik \emph{kanonik} tertentu. Himpunan-himpunan
kanonik ini pertama kali diperkenalkan oleh von Neumann:

\begin{defn}
Himpunan $A$ bersifat \emph{transitif} jika dan hanya jika $(\forall x \in
A)x \subseteq A$. Kemudian, $A$ merupakan suatu \emph{ordinal} jika dan hanya
jika $A$ bersifat transitif dan diurutkan dengan baik oleh $\in$.
\end{defn}
\noindent
Dalam pembahasan berikut, kita akan menggunakan huruf-huruf Yunani untuk
ordinal. Langsung dari definisi mengikuti bahwa, jika $\alpha$ merupakan
suatu ordinal, maka $\tuple{\alpha, \in_\alpha}$ merupakan suatu pengurutan
baik, dengan $\in_\alpha = \Setabs{\tuple{x, y} \in \alpha^2}{x \in y}$.
Jadi, dengan sedikit menyalahgunakan notasi, kita dapat mengatakan bahwa
$\alpha$ \emph{sendiri} merupakan suatu pengurutan baik.

Berikut beberapa ordinal pertama kita:
\[
	\emptyset, \{\emptyset\},
	\{\emptyset, \{\emptyset\}\},
	\{\emptyset, \{\emptyset\}, \{\emptyset, \{\emptyset\}\}\}, \ldots
\]
Anda akan memperhatikan bahwa ini merupakan beberapa ordinal pertama yang
kita jumpai dalam Aksioma Ketakhinggaan kita, yakni dalam definisi von
Neumann atas $\omega$ (lihat \olref[sth][z][infinity-again]{sec}). Hal ini
bukan kebetulan. Definisi von Neumann atas ordinal memperlakukan bilangan
asli sebagai ordinal, tetapi juga membolehkan ordinal transfinit.

Seperti biasa, kini kita dapat bertanya: \emph{apakah} semua ini merupakan
ordinal? Ataukah von Neumann sekadar memberi kita beberapa himpunan yang
dapat kita \emph{perlakukan} sebagai ordinal? Jenis pembahasan yang mungkin
dilakukan mengenai pertanyaan ini serupa dengan pembahasan kita dalam
\olref[sfr][rel][ref]{sec},
\olref[sfr][arith][ref]{sec},
\olref[sfr][infinite][dedekindsproof]{sec}, dan
\olref[sth][z][nat]{sec}, sehingga kita tidak akan berpanjang lebar mengenai
hal itu. Sebaliknya, dalam pembahasan berikut, kita cukup menggunakan
``ordinal'' untuk membicarakan ``ordinal von Neumann''.

\end{document}
